\section{Physical action, kinematics and liquids: proofs}
\label{sec:annex-p16-kr2008-action-and-causation}

These arguments use the exact physical/mathematical premises declared in
the reader. The survey's referrals to original programs and axiomatizations
are not thereby claimed as completely reconstructed publications.

\begin{annexitem}{annex:p16.phys-rc}{thm:p16-phys-rc}{Port composition, construction and uniqueness}
\begin{proof}
The resistor's incoming current at its \(b\) port is \((u-v)/R\).
Its incoming current at the shared \(a\) port is the negative. The
capacitor's incoming current at the shared \(b\) port is \(Cv'\).
The node sum is therefore \(-(u-v)/R+Cv'=0\), giving the stated equation.
Put \(k=1/(RC)>0\). Multiplying \(v'+kv=ku\) by \(e^{kt}\)
gives \((e^{kt}v)'=ke^{kt}u\). Integration and the initial value give
the displayed expression. Conversely continuity of \(u\) permits
differentiation by the fundamental theorem of calculus; the expression
has initial value \(v_0\) and satisfies the equation, so it constructs
a response. If two responses share the input and initial value, their
difference \(w\) satisfies \((e^{kt}w)'=0,w(0)=0\), and hence is zero.
Changing \(v_0\) while fixing \(u=0\) changes the exponential term.
Fixing \(v_0\) but choosing \(u=0\) versus \(u=1\) changes the
response by \(1-e^{-kt}\) at positive times. This verifies the boundary-
data qualification, without a claim about nonideal components.
\end{proof}
\end{annexitem}

\begin{annexitem}{annex:p16.phys-kinematics}{thm:p16-phys-kinematics}{All rigid invariants and both history directions}
\begin{proof}
The map \(x\mapsto Qx+b\) is a homeomorphism with inverse
\(y\mapsto Q^T(y-b)\). It preserves compactness and connectedness,
and commutes with closure and interior, so preserves regular closedness.
For poses \((Q,b),(Q',b')\) and a shape bounded by \(\|x\|\le M\),
the matching points in their two images differ by at most
\(M\|Q-Q'\|+\|b-b'\|\). Taking supremum distances in both
directions bounds their Hausdorff distance by that expression. Continuity
of pose therefore gives continuity of occupied regions.

A continuous path of feasible pose tuples assigns every source
\(Position(o)(t)\), keeps \(Shape(o)\) fixed, and defines
\(Place(o)(t)\) by the rigid image. The feasible-set condition enforces
noninterpenetration at every time, so the path constructs a history of the
declared model. Conversely every history has continuous pose tuples by
the model's position condition, and every tuple lies in \(\mathcal F\)
by noninterpenetration and the declared joints. It is thus a path with
the stated endpoints. Reparameterizing a positive finite time interval
to \([0,1]\) preserves these conditions. Concatenation and reversal
show that the reachable class is its path component. Mere connectedness
does not supply a path in an arbitrary topological subset; none was assumed.
\end{proof}
\end{annexitem}

\begin{annexitem}{annex:p16.phys-cells}{thm:p16-phys-cells}{Finite closed covers and exact path witnesses}
\begin{proof}
For a finite graph route, choose a configuration in each successive cell
intersection. Path connectedness joins the initial configuration to the
first witness, each witness to the next within its cell, and the last to
the target. Finitely concatenate and reparameterize those paths. Their
endpoints agree, so the result is continuous and remains in \(\mathcal F\).

Conversely let \(\gamma:[0,1]\to\mathcal F\) be a path. The sets
\(\gamma^{-1}(C_i)\) form a finite closed cover of \([0,1]\).
If cells containing its endpoints belonged to different graph components,
group these preimages by the component of the initial cell and its
complement. Both unions would be nonempty, closed and cover the interval.
They are disjoint: a shared time would put its configuration in two cells
whose intersection supplies an edge between the groups. This is a separation
of a connected interval, a contradiction. Any choices of endpoint cells
are themselves connected by their shared endpoint if more than one exists.

For failure without path-connected cells, take
\(\mathcal F=[-2,-1]\cup[1,2]\), cells \([-2,-1]\),
\(\{-1,1\}\), and \([1,2]\). Their intersection graph is a chain,
but any path between the endpoints would cross zero by the intermediate
value theorem. For failure with ambient-closure edges, take
\(\mathcal F=(-\infty,0)\cup(0,\infty)\) and its two component
cells. Their ambient closures meet at excluded zero, but no physical path
can cross it. The actual-intersection condition prevents that false edge.
\end{proof}
\end{annexitem}

\begin{annexitem}{annex:p16.phys-action-lift}{thm:p16-phys-action-lift}{All prescribed poses must belong to the same history}
\begin{proof}
A physical history accommodating the commanded objects gives a continuous
configuration path by the preceding construction, with its designated
coordinates equal to \(m(t)\). Thus it gives the stated lift. Conversely
such a lift constructs all the positions/places and a history performing
exactly those commanded motions, so the criterion is sufficient as well.
For the explicit constrained example, every negative \(x\) has frame
\((x,0)\), every positive \(x\) has \((x,1)\), and zero has both.
If a continuous lift existed, its second coordinate would be a continuous
function into \(\{0,1\}\). The intermediate value theorem makes it
constant on the interval. Its initial value is zero, but positive \(x\)
requires value one, a contradiction. Pointwise frame witnesses therefore
do not construct a feasible continuous action.
\end{proof}
\end{annexitem}

\begin{annexitem}{annex:p16.phys-nail}{thm:p16-phys-nail}{Force and torque conditions in both directions}
\begin{proof}
Vertical equilibrium requires the two upward friction forces to sum to
\(mg\). Each is at most \(\mu N\), so necessarily \(mg\le2\mu N\).
Conversely, if the inequality holds, choose both forces \(mg/2\).
They are admissible and balance weight. The equal opposing horizontal
normals balance horizontally. Symmetric friction forces have opposite
moments; equal-height opposing normals also have opposite moments, and
gravity through the central mass has none. All torque components therefore
balance. This constructs a static solution. Taking \(N=0\) violates
the necessary inequality, while any \(N\ge mg/(2\mu)\) works if
\(\mu>0\). The poses and initial zero velocities need not change
between these specified-force cases. No method of deriving the preload
from those poses was assumed.
\end{proof}
\end{annexitem}

\begin{annexitem}{annex:p16.phys-exit}{thm:p16-phys-exit}{The limiting set, dissipation and excluded escape faces}
\begin{proof}
Suppose a trajectory remains in \(K\) for all nonnegative time. Compactness
and continuity bound its energy below. Nonincrease makes the energy tend
to a finite value \(E_\infty\). The trajectory has a nonempty set
\(\Omega\) of limit points along sequences of times tending to infinity,
by compactness; \(\Omega\subseteq K\). Continuity gives energy
\(E_\infty\) at every point of \(\Omega\).
For \(z\in\Omega\), choose \(t_j\to\infty\) with trajectory
states tending to \(z\). Continuity and the flow law give
\(\Phi_\tau(z)=\lim_j\Phi_{t_j+\tau}(z_0)\in\Omega\) for
every fixed \(\tau\ge0\). Thus \(\Omega\) is forward invariant.
Energy on those orbits is constant, so its derivative \(D\) is zero
there. This contradicts the absence of a nonempty forward-invariant subset
of \(K\cap\{D=0\}\). A forever-trapped trajectory is impossible.
At a forbidden exit face, \(E\ge U>E(z_0)\) contradicts energy
nonincrease, proving the escape-face exclusion.

For the smooth damped model, differentiating gives
\[
 E'=v\cdot(-\nabla U-\gamma v)+\nabla U\cdot v
    =-\gamma\|v\|^2.
\]
An orbit remaining in its zero-derivative set has \(v=0\) throughout,
so \(q\) is constant; its acceleration equation then requires
\(\nabla U(q)=0\). If \(K\) contains no such equilibrium, the
zero set has no nonempty forward-invariant subset. Finally
\(x(t)=x(0)e^{-t}\) solves the counterexample, stays positive at every
finite time, and has energy derivative \(-2x(t)^2<0\). There is no
stationary point with \(x>0\), but it approaches the exit at infinite
time. In the compact closure, zero itself is a zero-dissipation invariant
set, exactly the condition the sufficient argument excludes.
\end{proof}
\end{annexitem}

\begin{annexitem}{annex:p16.phys-liquid-balance}{thm:p16-phys-liquid-balance}{Boundary transport, relative flux and a moving-region countercase}
\begin{proof}
For a smooth moving domain, the first-order volume change is the integral
of its boundary's outward normal velocity. To see this local transport
fact, use normal coordinates in a thin boundary tube: a patch swept
outward by normal displacement \(h\,u_n+o(h)\) adds its area times
that displacement, and an inward displacement subtracts it. The smooth
compact boundary gives uniform bounds for the tube Jacobian, which is
\(1+O(h)\); the remainder divided by \(h\) tends to zero.
Integrate the patches and let \(h\to0\). For the transverse intersection
\(L\cap r\), its boundary has a liquid-boundary part and a control-
boundary part. Their intersections have codimension two, and the small
swept neighborhoods there have volume \(O(h^2)\), so do not alter this
first derivative. Thus
\[
 V'=\int_{\partial L\cap r}v\cdot n_L\,dA
       +\int_{\partial r\cap L}w\cdot n_r\,dA.
\]
The divergence theorem on the intersection, with its two outward faces,
and incompressibility give
\[
 0=\int_{L\cap r}\nabla\cdot v\,dx
   =\int_{\partial L\cap r}v\cdot n_L\,dA
       +\int_{\partial r\cap L}v\cdot n_r\,dA.
\]
Subtract the second identity from the first to obtain
\(V'=-\int_{\partial r\cap L}(v-w)\cdot n_r\,dA=-F_{rel}\).
Reversing the surface orientation negates the integrand. With \(w=0\)
it is laboratory flux. For a material liquid domain alone, its boundary
velocity is \(v\); the same single-domain transport fact and divergence
theorem give volume derivative zero, without the transverse-intersection
premise needed for two different domains.

In the explicit interval example the overlap is \([0,t+1]\), of
length \(t+1\), throughout \(-1<t<1\). Its derivative is one.
The stationary liquid velocity gives zero laboratory flux. At the right
control endpoint the liquid is present, its outward orientation is positive,
and \(v-w=-1\); at the left endpoint no liquid is present. Relative
outflow is minus one, agreeing with the volume derivative. This verifies
why a moving boundary cannot simply be forgotten. Singular contacts,
coincident interfaces and topology changes were excluded, not proved by
the smooth tube argument.
\end{proof}
\end{annexitem}

\begin{annexitem}{annex:p16.phys-liquid-level}{thm:p16-phys-liquid-level}{Hydrostatic head and every permitted volume level}
\begin{proof}
The hydrostatic law gives \(\nabla(p+\rho gz)=0\). Along any piecewise
differentiable path in the fluid interior, the chain rule makes this
quantity constant. An open connected fluid domain admits polygonal paths:
the points reachable from a fixed one by a finite polygonal path form
a relatively open set, and their complement is relatively open by the
same small-ball continuation argument, so connectedness makes all points
reachable. Continuity extends the common pressure/head value to the free
surfaces. Consequently \(p_i+\rho gh_i=p_j+\rho gh_j\), and
rearrangement gives the displayed height relation. Its sign gives the
equal-pressure and higher-pressure conclusions.

The integral defining \(A\) is continuous and has \(A(h_0)=0\).
For \(h_1<h_2\), the stipulated strictly positive integral on their
interval gives \(A(h_2)-A(h_1)>0\), so it is strictly increasing.
The intermediate value theorem gives a level for every volume between
its two endpoint values; strict increase proves uniqueness, including
zero and the full below-opening capacity. No height below or at the
opening has capacity exceeding \(A(h_r)\), so a larger volume cannot
have a wholly contained hydrostatic realization under the premises.
This is not a theorem of inertial fluid motion: adding wetting, pressure
differences, inaccessible cavities or kinetic effects changes those premises.
\end{proof}
\end{annexitem}

\begin{annexitem}{annex:p16.phys-abstraction}{thm:p16-phys-abstraction}{Theory inclusion and a physically lost size constraint}
\begin{proof}
A model satisfying \(V\cup A\) satisfies \(V\), proving inclusion.
If all \(V\)-models satisfy a sentence, all the included \(U\)-models
do as well. For failure of reverse consistency, take \(V\) empty and
\(A=\{p,\neg p\}\): \(V\) has models but \(U\) does not.
This purely logical relation supplies neither an efficient solver nor
evidence that the added physical assumptions are true.

Put the corridor walls at horizontal coordinates zero and \(w\).
A disk center at \(x\) has interior extending \(r\) to each side.
Avoiding penetration requires \(x-r\ge0\) and \(x+r\le w\),
or \(r\le x\le w-r\). This interval is nonempty exactly when
\(w\ge2r\). At equality its unique horizontal center can translate
along the corridor with contact but no penetration. For a point, choose
horizontal coordinate \(w/2\) and translate longitudinally; this is
feasible whenever \(w>0\). The narrower-than-diameter corridor is thus
a concrete false-positive passage created by the point abstraction.
\end{proof}
\end{annexitem}

\begin{annexitem}{annex:p16.phys-braking}{thm:p16-phys-braking}{The full uncertainty range and both strict outcomes}
\begin{proof}
For \(v_0>0\) the stipulated speed is \(v_0-at\) until
\(t_s=v_0/a\), and then zero. Integrating speed gives
\(d=v_0t_s-at_s^2/2=v_0^2/(2a)\). At \(v_0=0\) it is already
stationary and the same distance formula gives zero. The formula is
nondecreasing in nonnegative speed and nonincreasing in positive
deceleration. Therefore its extrema are the two indicated corners, which
are permitted by the independent rectangle. Continuity on that connected
rectangle makes its image an interval containing both extrema, so it is
exactly the displayed full range.

Position increases to \(d\) and remains there. A candidate stops
strictly before the edge exactly when \(d<L\). It passes the edge
exactly when \(d>L\): its continuous increasing pre-stop position
then passes through \(L\) before stopping and becomes greater.
Taking universal or existential quantifiers over the closed attained
range gives all four bounds in the statement. Equality is neither
strictly-before nor past. If the edge lies strictly between the extrema,
the two corner realizations witness the opposite outcomes. Correlated
parameters may remove either corner; no assertion about their feasible
range follows from the rectangle calculation in that case.
\end{proof}
\end{annexitem}
